\section{The uniform bound and its threshold consequence}
\label{sec:completion}

We now choose the block size so that~\eqref{eq:captured-mass}
controls the low-degree part of the variance. The weighted Parseval
identity will control the remaining degrees.

\begin{proof}[Proof of Theorem~\ref{thm:main}]
Fix $r\ge3$ and write $\alpha=r/(r-1)$. Choose an integer $N_r>r$
such that, for every $n\ge N_r$ and $m=\lfloor\sqrt n\rfloor$,
\begin{equation}\label{eq:large-n}
 \begin{gathered}
 m\ge r+1,\qquad \log m\ge\tfrac13\log n,\\
 r(\log n)^\alpha n^{-1/2}\le\tfrac12,\qquad
 m^{-1/4}(\log n)^\alpha\le1.
 \end{gathered}
\end{equation}
Such an $N_r$ exists: $m$ is asymptotic to $\sqrt n$, and any
fixed power of $\log n$ grows more slowly than a positive power of
$n$. In particular, the choice is independent of $p$ and $f$.

Assume first that $n\ge N_r$, and set
\[
 \sigma=\sqrt{p(1-p)},\qquad
 k=\frac{\sigma\log m}{64},\qquad L=(k/r)^\alpha.
\]
Both cutoffs are positive. Since $k/r\le\log n$ and
\[
 \frac{m-1}{n-1}\le\frac{\sqrt n-1}{n-1}
 =\frac1{\sqrt n+1}\le n^{-1/2},
\]
the third condition in~\eqref{eq:large-n} implies
\eqref{eq:capture-condition}. Thus~\eqref{eq:captured-mass} and
the fourth condition in~\eqref{eq:large-n} give
\begin{equation}\label{eq:low-final}
 \sum_{1\le|S|\le L}\widehat f(S)^2
 \le\frac{2r}{(\log n)^\alpha}I_p(f).
\end{equation}

For the remaining degrees,~\eqref{eq:parseval} gives
\[
 \sum_{|S|>L}\widehat f(S)^2
 \le\frac1L\sum_S|S|\widehat f(S)^2
 =\frac{\sigma^2}{L}I_p(f).
\]
The decisive bias dependence is
\begin{equation}\label{eq:bias-cancellation}
 \frac{\sigma^2}{L}
 =\frac{(64r)^\alpha\sigma^{2-\alpha}}{(\log m)^\alpha}
 \le\frac{(192r)^\alpha}{(\log n)^\alpha}.
\end{equation}
Indeed $2-\alpha=(r-2)/(r-1)>0$, $\sigma\le1$, and
$\log m\ge(\log n)/3$. This estimate remains valid at arbitrarily
small positive $\sigma$. If $L<1$, the low-degree sum is empty and
the same tail estimate bounds the whole variance.

Adding the two parts proves~\eqref{eq:main} for $n\ge N_r$ with
constant $2r+(192r)^\alpha$. For $r\le n<N_r$,
\eqref{eq:parseval} gives $\Var_p(f)\le\sigma^2 I_p(f)\le I_p(f)/4$.
Therefore a valid choice for all $n$ and $p$ is
\[
 C_r=\max\left\{2r+(192r)^\alpha,\,
                       \frac{(\log N_r)^\alpha}{4}\right\}.
\]
\end{proof}

We finish with the threshold consequence stated in the introduction.
Recall that $q(p)=\E_p f$ and
$p_a=\inf\{p\in[0,1]:q(p)\ge a\}$.

\begin{corollary}\label{cor:threshold}
For a nontrivial increasing property of $r$-uniform hypergraphs on
$n\ge r$ vertices, where $r\ge3$, and every $0<\varepsilon<1/2$,
\[
 p_{1-\varepsilon}-p_\varepsilon
 \le\frac{2C_r}{(\log n)^{r/(r-1)}}
                  \log\frac{1-\varepsilon}{\varepsilon}.
\]
\end{corollary}

\begin{proof}
The Margulis--Russo identity relates pivotal influences to the
probability derivative; see~\cite{Margulis74},
\cite[Lemma~3]{Russo81}, and \cite[Lemma~2.2]{FK96}.
Here it follows directly by giving each coordinate
its own bias: the expectation is affine in each bias separately,
and its partial derivative in coordinate $e$ is
$\E_p[f(x_{e\leftarrow1})-f(x_{e\leftarrow0})]$.
Monotonicity identifies this with $I_{e,p}(f)$. Differentiating
along the common bias gives $q'(p)=I_p(f)$.

Nontrivial monotonicity gives $q(0)=0$ and $q(1)=1$. Some pair of
configurations differing in one coordinate has different values,
so its pivotal event has positive probability for every $0<p<1$.
Hence $q'(p)>0$ there, and the stated quantiles exist uniquely.
Theorem~\ref{thm:main} implies
\[
 \frac{q'(p)}{q(p)(1-q(p))}
 \ge\frac{(\log n)^{r/(r-1)}}{C_r}.
\]
Integrating from $p_\varepsilon$ to $p_{1-\varepsilon}$, using
the antiderivative $\log(q/(1-q))$, proves the claim.
\end{proof}
